\documentclass[11pt]{article}

\usepackage{epsfig}
\usepackage{amsfonts}
\usepackage{amssymb}
\usepackage{amstext}
\usepackage{amsmath}
\usepackage{xspace}
\usepackage{theorem}
\usepackage{times}
\usepackage{graphicx}

%\usepackage{layout}% if you want to see the layout parameters
                     % and now use \layout command in the body

% This is the stuff for normal spacing
\makeatletter
 \setlength{\textwidth}{6.5in}
 \setlength{\oddsidemargin}{0in}
 \setlength{\evensidemargin}{0in}
 \setlength{\topmargin}{0.25in}
 \setlength{\textheight}{8.25in}
 \setlength{\headheight}{0pt}
 \setlength{\headsep}{0pt}
 \setlength{\marginparwidth}{59pt}

 \setlength{\parindent}{0pt}
 \setlength{\parskip}{5pt plus 1pt}
 \setlength{\theorempreskipamount}{5pt plus 1pt}
 \setlength{\theorempostskipamount}{0pt}
 \setlength{\abovedisplayskip}{8pt plus 3pt minus 6pt}

 \renewcommand{\section}{\@startsection{section}{1}{0mm}%
                                   {2ex plus -1ex minus -.2ex}%
                                   {1.3ex plus .2ex}%
                                   {\normalfont\Large\bfseries}}%
 \renewcommand{\subsection}{\@startsection{subsection}{2}{0mm}%
                                     {1ex plus -1ex minus -.2ex}%
                                     {1ex plus .2ex}%
                                     {\normalfont\large\bfseries}}%
 \renewcommand{\subsubsection}{\@startsection{subsubsection}{3}{0mm}%
                                     {1ex plus -1ex minus -.2ex}%
                                     {1ex plus .2ex}%
                                     {\normalfont\normalsize\bfseries}}
 \renewcommand{\paragraph}{\@startsection{paragraph}{4}{0mm}%
                                    {1ex \@plus1ex \@minus.2ex}%
                                    {-1em}%
                                    {\normalfont\normalsize\bfseries}}
 \renewcommand{\subparagraph}{\@startsection{subparagraph}{5}{\parindent}%
                                       {2.0ex \@plus1ex \@minus .2ex}%
                                       {-1em}%
                                      {\normalfont\normalsize\bfseries}}
\makeatother

\newenvironment{proof}{{\bf Proof:  }}{\hfill\rule{2mm}{2mm}}
\newenvironment{proofof}[1]{{\bf Proof of #1:  }}{\hfill\rule{2mm}{2mm}}
\newenvironment{proofofnobox}[1]{{\bf#1:  }}{}
\newenvironment{example}{{\bf Example:  }}{\hfill\rule{2mm}{2mm}}
\renewcommand{\thesection}{\lecnum.\arabic{section}}

\renewcommand{\theequation}{\thesection.\arabic{equation}}
\renewcommand{\thefigure}{\thesection.\arabic{figure}}

\newtheorem{fact}{Fact}[section]
\newtheorem{lemma}[fact]{Lemma}
\newtheorem{theorem}[fact]{Theorem}
\newtheorem{definition}[fact]{Definition}
\newtheorem{corollary}[fact]{Corollary}
\newtheorem{proposition}[fact]{Proposition}
\newtheorem{claim}[fact]{Claim}
\newtheorem{exercise}[fact]{Exercise}

% math notations
\newcommand{\R}{\ensuremath{\mathbb R}}
\newcommand{\Z}{\ensuremath{\mathbb Z}}
\newcommand{\N}{\ensuremath{\mathbb N}}
\newcommand{\F}{\ensuremath{\mathcal F}}
\newcommand{\SymGrp}{\ensuremath{\mathfrak S}}

\newcommand{\size}[1]{\ensuremath{\left|#1\right|}}
\newcommand{\ceil}[1]{\ensuremath{\left\lceil#1\right\rceil}}
\newcommand{\floor}[1]{\ensuremath{\left\lfloor#1\right\rfloor}}
\newcommand{\poly}{\operatorname{poly}}
\newcommand{\polylog}{\operatorname{polylog}}

% asymptotic notations
\newcommand{\Oh}[1]{{\mathcal O}\left({#1}\right)}
\newcommand{\LOh}[1]{{\mathcal O}\left({#1}\right.}
\newcommand{\ROh}[1]{{\mathcal O}\left.{#1}\right)}
\newcommand{\oh}[1]{{o}\left({#1}\right)}
\newcommand{\Om}[1]{{\Omega}\left({#1}\right)}
\newcommand{\om}[1]{{\omega}\left({#1}\right)}
\newcommand{\Th}[1]{{\Theta}\left({#1}\right)}


% pseudocode notations
\newcommand{\xif}{{\bf{\em{if~}}}}
\newcommand{\xthen}{{\bf{\em{then~}}}}
\newcommand{\xelse}{{\bf{\em{else~}}}}
\newcommand{\xelseif}{{\bf{\em{elif~}}}}
\newcommand{\xfi}{{\bf{\em{fi~}}}}
\newcommand{\xcase}{{\bf{\em{case~}}}}
\newcommand{\xendcase}{{\bf{\em{endcase~}}}}
\newcommand{\xfor}{{\bf{\em{for~}}}}
\newcommand{\xto}{{\bf{\em{to~}}}}
\newcommand{\xby}{{\bf{\em{by~}}}}
\newcommand{\xdownto}{{\bf{\em{downto~}}}}
\newcommand{\xdo}{{\bf{\em{do~}}}}
\newcommand{\xrof}{{\bf{\em{rof~}}}}
\newcommand{\xwhile}{{\bf{\em{while~}}}}
\newcommand{\xendwhile}{{\bf{\em{endwhile~}}}}
\newcommand{\xand}{{\bf{\em{and~}}}}
\newcommand{\xor}{{\bf{\em{or~}}}}
\newcommand{\xerror}{{\bf{\em{error~}}}}
\newcommand{\xreturn}{{\bf{\em{return~}}}}
\newcommand{\xparallel}{{\bf{\em{parallel~}}}}
\newcommand{\T}{\hspace{0.5cm}}
\newcommand{\m}{\mathcal}

\def\sland{~\land~}
\def\slor{~\lor~}
\def\sRightarrow{~\Rightarrow~}

\def\comment#1{\hfill{$\left\{\textrm{{\em{#1}}}\right\}$}}
\def\lcomment#1{\hfill{$\left\{\textrm{{\em{#1}}}\right.$}}
\def\rcomment#1{\hfill{$\left.\textrm{{\em{#1}}}\right\}$}}
\def\fcomment#1{\hfill{$\textrm{{\em{#1}}}$}}
\def\func#1{\textrm{\bf{\sc{#1}}}}
\def\funcbf#1{\textrm{\textbf{\textsc{#1}}}}

\newcommand{\hide}[1]{}

\newcommand{\prob}[1]{\ensuremath{\text{{\bf Pr}$\left[#1\right]$}}}
\newcommand{\expct}[1]{\ensuremath{\text{{\bf E}$\left[#1\right]$}}}
\newcommand{\Event}{{\mathcal E}}

\newcommand{\mnote}[1]{\normalmarginpar \marginpar{\tiny #1}}

\makeatletter
   \newcommand\figcaption{\def\@captype{figure}\caption}
   \newcommand\tabcaption{\def\@captype{table}\caption}
\makeatother


%%%%%%%%%%%%%%%%%%%%%%%%%%%%%%%%%%%%%%%%%%%%%%%%%%%%%%%%%%%%%%%%%%%%%%%%%%%
% Document begins here %%%%%%%%%%%%%%%%%%%%%%%%%%%%%%%%%%%%%%%%%%%%%%%%%%%%
%%%%%%%%%%%%%%%%%%%%%%%%%%%%%%%%%%%%%%%%%%%%%%%%%%%%%%%%%%%%%%%%%%%%%%%%%%%

\newcommand{\headings}[4]{
{\bf CSE548: Analysis of Algorithms, Fall 2025} \hfill {{\bf Lecturer:} #1}\\
{{\bf Topic:} #2} \hfill {{\bf Date:} #3} \\
{{\bf Scribe:} #4}\\
\rule[0.1in]{\textwidth}{0.025in}
%\thispagestyle{empty}
}

\begin{document}

\headings{Name of the Lecturer}{Lecture Topic}{Lecture Date}{Name of the Scribe}
\newcommand{\lecnum}{5}  % Lecture Number

\section{This is the First Section of Lecture \lecnum}

The text goes here.

Here is how to include a theorem with a proof.

\begin{theorem}
 When executed with $p$ processors,
 multithreaded IGEP performs $T_{1} = \Oh{n^3}$ work and terminates
 in ${\frac{T_{1}}{p}} + T_{\infty} = \Oh{\frac{n^3}{p}} + n\log^{2}{n}$
 parallel steps on an $n \times n$ input matrix.
\label{thm:igep}
\end{theorem}

\begin{proof}
This is a proof of Theorem \ref{thm:igep}.
\end{proof}


This is how to include {\tt eps} files that can be generated using {\tt xfig}.

\begin{figure}[h]
\begin{center}
%\epsfig{file = figs/testfig.eps, height = 1in}
\figcaption{Include EPS files.} 
\end{center}
\end{figure}


You can include {\tt{pdf}}, {\tt{jpg}}, {\tt{png}} files as follows.

\begin{figure*}[h!]
\centering
%\includegraphics[width=\textwidth]{figs/testfig.pdf}
\figcaption{Include PDF/JPG/PNG files.} 
\label{fig:example-figure}
\end{figure*}


Figure \ref{fig:ApproxInterfaceArea} shows how to write pseudocode.

 \begin{figure*}[t!] 
%
 \begin{minipage}{\textwidth}
 \begin{center}
 \framebox{
 \begin{minipage}{\textwidth}
 {\scriptsize
 \vspace{-0.2cm}
 \medskip\noindent\func{CollectInterfacePatches$(~x, ~y, ~\m{PG}~)$}

 \vspace{0.1cm}
 \noindent
(Inputs are two octree nodes $x \in \m{T}_{S_A}$ and $y \in \m{T}_{S_B}$.
The function collects all interface patches of $S_A$ and $S_B$
in a packing grid data structure $\m{PG}$.
%
By $\func{child}(x)$ (resp. $\func{child}(y)$) we denote the set of non-empty octree nodes obtained by
subdividing node $x$ (resp. $y$). We denote by $r_{x, y}$ the distance between the
centers of $x$ and $y$.)

 \vspace{0.2cm}
 \noindent
 \begin{enumerate}

 \item \xif $r_{x,y} - (r_{x} + r_{y}) \geq \mu$ \xthen \xreturn
%
             \comment{too far: no interface between $\m{N}_x$ and $\m{N}_y$}

 \item \xelse \xif $\func{leaf}(x)$ $\wedge$ $\func{leaf}(y)$ \xthen 
%
             \comment{collect interface patches, if any}

             \T\T \xfor each $\pi \in \m{N}_x$ \xdo

              \T\T\T \xif $\left( \pi \notin \m{PG} \right) \wedge \exists_{p \in \m{PG}_{y}}{\left( dist( c_{\pi}, p ) < \mu \right)} $ \xthen
 %
                        insert $\pi$ into $\m{PG}$
%
             \comment{collect interface patches of $S_A$}
 
              \T\T \xfor each $\pi \in \m{N}_y$ \xdo
 
              \T\T\T \xif $\left( \pi \notin \m{PG} \right) \wedge \exists_{p \in \m{PG}_{x}}{\left( dist( c_{\pi}, p ) < \mu \right)} $ \xthen
% 
                        insert $\pi$ into $\m{PG}$
%
             \comment{collect interface patches of $S_B$}

 \item \T ~\xelse \xif $\func{leaf}(x)$ \xthen 

             \T\T\T \xfor each $cy \in \func{child}(y)$ \xdo $\func{CollectInterfacePatches}(~x, ~cy, ~\m{PG}~)$
%
             \comment{recursive collection}

 \item \T\T ~\xelse \xif $\func{leaf}(y)$ \xthen

             \T\T\T\T \xfor each $cx \in \func{child}(x)$ \xdo $\func{CollectInterfacePatches}(~cx, ~y, ~\m{PG}~)$
%
             \comment{recursive collection}

 \item \T\T\T ~\xelse 
%
              \xfor all $cx \in \func{child}(x)$ \xand $cy \in \func{child}(y)$ \xdo 

               \T\T\T\T\T $\func{CollectInterfacePatches}(~cx, ~cy, ~\m{PG}~)$
%
             \comment{recursive collection}

 \end{enumerate}
 \medskip\noindent{\func{CollectInterfacePatches Ends}}
 }
 \end{minipage}
 }
 \end{center}
 \end{minipage}

 \begin{minipage}{\textwidth}
 \begin{center}
 \framebox{
 \begin{minipage}{\textwidth}
 {\scriptsize
 \vspace{-0.2cm}
 \medskip\noindent\func{ApproxInterfaceArea$(~A, ~B~)$}

 \vspace{0.1cm}
 \noindent
(Inputs are two molecules $A$ and $B$. We assume that for each 
molecule $P \in \{ A, B \}$, the octree $\m{T}_{S_{P}}$
and the packing grid $\m{PG}_{S_{P}}$
based on the patches on its surface $S_{P}$
have already been constructed.
The function evaluates and returns $InterfaceArea( A, B )$.)

 \vspace{0.2cm}
 \noindent
 \begin{enumerate}

 \item initialize an empty packing grid $\m{PG}$
 %
 \comment{for collecting interface patches without any duplicates}

 \item $\func{CollectInterfacePatches}(~\func{root}(\m{T}_{S_A}),~\func{root}(\m{T}_{S_B}), ~\m{PG}~)$
 %
 \comment{collect the interface patches in $\m{PG}$}

 \item \xreturn $\sum_{\pi \in \m{PG}}{a_{\pi}}$
 %
 \comment{sum up the patch areas and return}

 \end{enumerate}
 \medskip\noindent{\func{ApproxInterfaceArea Ends}}
 }
 \end{minipage}
 }
 \end{center}
 \end{minipage}

 \figcaption{\small {\sc(Example Pseudocode)} Identify the interface patches on $S_A$ and $S_B$ 
 and return the sum of their areas. 
 The function \func{CollectInterfacePatches$(~x, ~y, ~\m{PG}~)$}
 collects in $\m{PG}$ (without duplicates) each patch from $\m{N}_x$ 
 that is within distance $\mu$ from at least one patch in $\m{N}_y$ 
 (center to center distance), and vice versa. 
 The initial call is $\func{CollectInterfacePatches}(~\func{root}(\m{T}_{S_A}),~\func{root}(\m{T}_{S_B}), ~\m{PG}~)$
 for computing $InterfaceArea( A, B )$.}
 \label{fig:ApproxInterfaceArea}
 \end{figure*}

\end{document}
